\documentclass[11pt,a4paper]{article}

\usepackage[T1]{fontenc}
\usepackage[utf8]{inputenc}
\usepackage[margin=25mm]{geometry}
\usepackage{enumitem}
\usepackage[protrusion=true,expansion=false]{microtype}
\usepackage[hidelinks]{hyperref}
\usepackage{amsmath,amssymb}

\title{Calibration Dataset Adequacy Framework\\Working Research Plan}
\author{}
\date{\today}

\begin{document}

\maketitle

\section{Objective}

The objective is to develop a practical framework for assessing whether a
planned experiment or a collected dataset is adequate for calibrating a sensor
for a declared use.

Adequacy is not absolute. It depends on what the sensor must measure, the
intended operating range and conditions, the calibration method, the required
performance, and the situations in which the calibrated sensor is expected to
work.

\section{Central Research Question}

\begin{quote}
Given a calibration objective, how can we determine whether an experimental
protocol or dataset provides enough appropriate evidence to develop and
evaluate a viable sensor calibration?
\end{quote}

\section{Working Definition of Adequacy}

A dataset cannot be considered adequate in isolation. Its adequacy must be
assessed relative to a declared calibration task:
\begin{equation}
    \mathcal{T}
    =
    \left(
        \mathbf{s},
        \mathbf{q},
        \Omega,
        \mathcal{C},
        \mathcal{M},
        \mathcal{P},
        \mathcal{G}
    \right),
\end{equation}
where:
\begin{itemize}
    \item $\mathbf{s}$ is the vector of raw sensor signals used as calibration
        inputs;
    \item $\mathbf{q}$ is the reference measurand that the calibration must
        estimate;
    \item $\Omega$ is the intended operating domain;
    \item $\mathcal{C}$ is the set of relevant operating conditions;
    \item $\mathcal{M}$ is the calibration model or candidate model family;
    \item $\mathcal{P}$ is the required performance and uncertainty; and
    \item $\mathcal{G}$ is the intended generalization claim.
\end{itemize}

For an acquired dataset $D$, adequacy is therefore conditional:
\begin{equation}
    \operatorname{Adequacy}\left(D \mid \mathcal{T}\right).
\end{equation}

\section{Trixial Force Sensor (under rigid mounting)}

\subsection{Adequacy Definition}

For the first calibration stage, the sensor input and reference output are
\begin{equation}
    \mathbf{s}
    =
    \begin{bmatrix}
        c_1 & c_2 & c_3
    \end{bmatrix}^{\mathsf{T}}
    \in \mathcal{S} \subseteq \mathbb{R}^{3},
    \qquad
    \mathbf{F}
    =
    \begin{bmatrix}
        F_x & F_y & F_z
    \end{bmatrix}^{\mathsf{T}}
    \in \mathbb{R}^{3}.
\end{equation}

The calibration model estimates the reference force according to
\begin{equation}
    \widehat{\mathbf{F}} = f_{\boldsymbol{\theta}}(\mathbf{s}),
\end{equation}
where $f_{\boldsymbol{\theta}}$ is a candidate calibration model with
parameters $\boldsymbol{\theta}$.

\subsubsection{Intended Force Domain}

The intended three-dimensional force domain is the set of force vectors for
which the calibration will be claimed:
\begin{equation}
    \Omega
    =
    \left\{
        \mathbf{F} \in \mathbb{R}^{3}
        \;\middle|\;
        \mathbf{F}^{-}
        \preceq
        \mathbf{F}
        \preceq
        \mathbf{F}^{+},
        \quad
        h_k(\mathbf{F}) \leq 0,
        \quad k=1,\ldots,K
    \right\}.
\end{equation}

Here,
\begin{equation}
    \mathbf{F}^{-}
    =
    \begin{bmatrix}
        F_x^{-} & F_y^{-} & F_z^{-}
    \end{bmatrix}^{\mathsf{T}},
    \qquad
    \mathbf{F}^{+}
    =
    \begin{bmatrix}
        F_x^{+} & F_y^{+} & F_z^{+}
    \end{bmatrix}^{\mathsf{T}},
\end{equation}
are the lower and upper force limits, and $\preceq$ denotes component-wise
inequality. The functions $h_k$ represent any coupled physical or
application-specific restrictions on valid force combinations. If no such
restrictions are declared, $K=0$ and $\Omega$ is the rectangular domain defined
only by the force limits.

\subsubsection{Mounting and Operating Conditions}

Let
\begin{equation}
    \mathbf{z}
    =
    \begin{bmatrix}
        z_1 & z_2 & \cdots & z_r
    \end{bmatrix}^{\mathsf{T}}
\end{equation}
represent the operating-condition variables that may affect the sensor
response. Their declared domain is
\begin{equation}
    \mathcal{C}
    =
    \mathcal{C}_1
    \times
    \mathcal{C}_2
    \times
    \cdots
    \times
    \mathcal{C}_r.
\end{equation}

For the first triaxial calibration stage, a possible parameterization is
\begin{equation}
    \mathbf{z}
    =
    \left(
        m,
        T,
        \dot{\mathbf{F}},
        h
    \right),
\end{equation}
where $m$ is the mounting configuration, $T$ is temperature,
$\dot{\mathbf{F}}$ is loading rate, and $h$ represents relevant loading
history. The corresponding condition set may be written as
\begin{equation}
    \mathcal{C}
    =
    \left\{m_{\mathrm{rigid}}\right\}
    \times
    \left[T^{-},T^{+}\right]
    \times
    \mathcal{V}
    \times
    \mathcal{H},
\end{equation}
where $\mathcal{V}$ is the allowed set of loading rates and $\mathcal{H}$ is
the allowed set of loading histories. A condition variable should be retained
only if it is relevant to the intended calibration claim.

\subsubsection{Candidate Calibration Models}

The candidate model family is
\begin{equation}
    \mathcal{M}
    =
    \bigcup_{k \in \mathcal{K}} \mathcal{M}_k,
    \qquad
    \mathcal{M}_k
    =
    \left\{
        f_{k,\boldsymbol{\theta}}
        \;\middle|\;
        \boldsymbol{\theta} \in \Theta_k
    \right\},
\end{equation}
where $\mathcal{M}_k$ is a model class and $\Theta_k$ is its parameter space.

The initial memoryless affine model class is
\begin{equation}
    \mathcal{M}_{\mathrm{affine}}
    =
    \left\{
        \mathbf{s}
        \mapsto
        \mathbf{A}\mathbf{s}+\mathbf{b}
        \;\middle|\;
        \mathbf{A}\in\mathbb{R}^{3\times3},
        \;
        \mathbf{b}\in\mathbb{R}^{3}
    \right\}.
\end{equation}

A full quadratic model can be defined using
\begin{equation}
    \boldsymbol{\phi}_2(\mathbf{s})
    =
    \begin{bmatrix}
        1 &
        c_1 & c_2 & c_3 &
        c_1^2 & c_2^2 & c_3^2 &
        c_1c_2 & c_1c_3 & c_2c_3
    \end{bmatrix}^{\mathsf{T}},
\end{equation}
giving
\begin{equation}
    \mathcal{M}_{\mathrm{quadratic}}
    =
    \left\{
        \mathbf{s}
        \mapsto
        \mathbf{B}\boldsymbol{\phi}_2(\mathbf{s})
        \;\middle|\;
        \mathbf{B}\in\mathbb{R}^{3\times10}
    \right\}.
\end{equation}

More complex nonlinear or machine-learning model classes may be introduced
later only when simpler model classes are shown to be inadequate.

\subsubsection{Required Performance and Uncertainty}

For validation sample $i$, define the force-estimation error as
\begin{equation}
    \mathbf{e}_i
    =
    \widehat{\mathbf{F}}_i
    -
    \mathbf{F}^{\mathrm{ref}}_i.
\end{equation}

For axis $j\in\{x,y,z\}$, the root-mean-square error is
\begin{equation}
    \operatorname{RMSE}_j
    =
    \sqrt{
        \frac{1}{N_{\mathrm{val}}}
        \sum_{i=1}^{N_{\mathrm{val}}}
        e_{i,j}^{2}
    }.
\end{equation}

The performance requirements are represented by a set of metric--threshold
pairs:
\begin{equation}
    \mathcal{P}
    =
    \left\{
        m_{\ell}
        \leq
        \tau_{\ell}
        \;:\;
        \ell=1,\ldots,L
    \right\},
\end{equation}
where $m_{\ell}$ is a declared validation metric and $\tau_{\ell}$ is its
acceptance threshold.

For example, axis-specific requirements may include
\begin{equation}
    \operatorname{RMSE}_j \leq \varepsilon^{\mathrm{RMSE}}_j,
    \qquad
    Q_{1-\alpha}\!\left(\lvert e_j\rvert\right)
    \leq
    \varepsilon^{\mathrm{quantile}}_j,
    \qquad
    U_j \leq U_j^{\max},
\end{equation}
where $Q_{1-\alpha}$ is the $(1-\alpha)$ error quantile and $U_j$ is the
declared calibration or prediction uncertainty for axis $j$. The metrics,
thresholds, and confidence level must be selected from the application
requirements rather than from the observed calibration results.

\subsubsection{Generalization Claim}

Let $\mathcal{R}$ denote the set of experimental runs, $\mathcal{C}$ the set of
operating conditions, and $\mathcal{U}$ the set of physical sensor units. The
generalization target is defined as
\begin{equation}
    \mathcal{G}
    \subseteq
    \mathcal{R}
    \times
    \mathcal{C}
    \times
    \mathcal{U}.
\end{equation}

For a calibration intended to generalize to new runs of the same individual
sensor under the declared conditions,
\begin{equation}
    \mathcal{G}_{\mathrm{stage\,1}}
    =
    \mathcal{R}_{\mathrm{new}}
    \times
    \mathcal{C}
    \times
    \left\{u_0\right\},
\end{equation}
where $u_0$ is the calibrated sensor unit.

If the calibration is later required to generalize to previously unseen sensor
units, the claim becomes
\begin{equation}
    \mathcal{G}_{\mathrm{multi\text{-}unit}}
    =
    \mathcal{R}_{\mathrm{new}}
    \times
    \mathcal{C}
    \times
    \mathcal{U}_{\mathrm{new}}.
\end{equation}

The validation partition must match the generalization claim. For example,
generalization to new runs requires holding out complete runs, while
generalization to new sensor units requires holding out complete sensor units.


\end{document}

Where s is the readings of the three coils (c1, c2,c3) and q represent the three axial forces (fx, Fy, Fz).